\usepackage{graphicx}
\usepackage{geometry}
\geometry{
    a4paper,
    top = 2cm,
    left = 2cm,
    right = 2cm,
    bottom = 3cm
}
\usepackage[utf8]{inputenc}
\usepackage[spanish]{babel} % Renombra en automático los elementos a español
\usepackage{float}  % Etiqueta 'H' (mayus) en figuras (imágenes)
\usepackage{array}  % Padding en tablas
\usepackage{url}    % Etiqueta url en .bib
\usepackage{booktabs}   % Manejo de caudros (tablas)
\usepackage[T1]{fontenc}    % Codificación de fuentes
\usepackage{csquotes}       % Paquete para comillas tipográficas
\usepackage{comment}    % Comentar bloques de líneas mediante begin-end
\usepackage{multicol}   % Para poner contenido a doble columna
\usepackage{enumitem}   % Para utilizar [noitemsep, topsep=0pt] en itemize
\usepackage{amsmath}
\usepackage[hidelinks]{hyperref}    % Habilitar hipervínculos sin cambiar colores
\usepackage{mathabx} % Para testar en math mode
\usepackage{makecell}
\usepackage{fancyhdr}
\usepackage[backend=biber, style=apa]{biblatex}
\addbibresource{Referencias/Ref.bib}

% ... código anterior ...
\addto\captionsspanish{\renewcommand{\tablename}{Tabla}}
\addto\captionsspanish{\def\tablenamepunct{. }}
% \addto\captionsspanish{\def\@captionhead#1#2{#1 #2.}} % Línea comentada

%--------------- Configuraciones para tablas ---------------
\usepackage{adjustbox}
\usepackage[table,xcdraw]{xcolor}
\definecolor{cellColor}{rgb}{0.788, 0.855, 0.972}    % RGB: 201, 218, 248 (Azul ocuro, Texto 2, Claro 80%)
% \cellcolor{cellColor} -> ELIMINADA
\usepackage{multirow}
% ----------------------------------------------------------